Un alambre de longitud $\pi\,\mathrm{cm}$ tiene difusividad térmica $1\,\mathrm{cm^2/s}$, superficie lateral aislada y extremos mantenidos a $0\,{}^\circ\mathrm C$. Inicialmente su temperatura interior es $25\,{}^\circ\mathrm C$. Una corriente eléctrica produce calor a razón de $5\,{}^\circ\mathrm C/\mathrm s$ en cada punto. Formule el modelo y determine $u(x,t)$ para $0<x<\pi$ y $t>0$.

Se usa conducción unidimensional homogénea, con difusividad constante y sin pérdidas laterales. La condición inicial uniforme corresponde al interior; las temperaturas de los extremos se imponen para $t>0$.
